\part{Pathophysiology and Biological Mechanisms}
\label{part:pathophysiology}

This part explores the biological underpinnings of ME/CFS, from well-established phenomena to emerging theories. We examine:

\begin{itemize}
\item \textbf{Known phenomena}: Mechanisms with strong research support
\item \textbf{Suspected phenomena}: Plausible mechanisms with preliminary evidence
\item \textbf{Related phenomena}: Connections to other conditions (allergies, autoimmune diseases, etc.)
\item \textbf{Biochemical processes}: Detailed molecular and cellular mechanisms
\item \textbf{Causal hierarchy}: Which mechanisms are root causes versus amplifiers versus downstream consequences, and why this distinction matters for treatment prioritization
\end{itemize}

Understanding these mechanisms is crucial for developing targeted treatments and explaining the diverse symptomatology of ME/CFS. The part concludes with a capstone synthesis (Chapter~\ref{ch:causal-hierarchy}) that classifies all mechanisms into a causal hierarchy, distinguishing those capable of triggering ME/CFS from those that merely amplify or result from the disease process.

